\ch{Chapter 2: Literature Review}\label{chap:chapter_2}



The literature on Digital Engineering tells a story of demonstrated capability confined within disciplinary walls. Model-Based Systems Engineering, digital threads, digital twins, and Product Lifecycle Management have matured through decades of application in defense and aerospace, accumulating institutional endorsement from the Department of War, NASA, and INCOSE alongside sparse but consistent empirical evidence of return on investment. Yet this maturation has occurred almost entirely within the systems engineering discipline. Enterprise IT infrastructure management and Information Assurance confront structurally analogous challenges of documentation accuracy, configuration visibility, compliance traceability, and change coordination. Nevertheless, these domains have developed independently, producing their own frameworks, tools, and patterns of persistent failure. The result is a pronounced research gap. Two efforts approach the problem: a preliminary reference model by \textcite{Bonar_Hastings_2024} and a conceptual framework for DevSecOps security assurance from the Carnegie Mellon Software Engineering Institute~\parencite{CMU_SEI_DevSecOps_2023}. Beyond these, no peer-reviewed research identified in this review addresses applying Digital Engineering methods to enterprise IT infrastructure management and Information Assurance programs as an integrated enterprise discipline. That scope remains distinct from Model-Based Systems Engineering applied to individual engineered systems or to single compliance controls.

This chapter examines why that gap persists and why it matters. This structured, analytical review spans disciplinary evolution, MBSE value evidence, adoption dynamics, cross-disciplinary transfer barriers, digital twin security applications, compliance automation, IT service management failures, and enterprise visibility challenges. Rather than cataloging what exists, the review foregrounds the debates, tensions, and unresolved questions that position this dissertation's investigation within the academic conversation. Four interrelated arguments emerge. First, systems engineering, Information Assurance, cybersecurity, and IT service management evolved along independent academic and professional trajectories, creating disciplinary silos that explain the absence of cross-disciplinary research. Second, the evidence base for MBSE value, while growing, remains disproportionately reliant on perceived rather than measured outcomes, creating economic uncertainty that discourages adoption beyond established domains. Third, adoption research shows that perceptions, rather than objective technical merit, drive adoption decisions, yet no empirical data exist on how IT and Information Assurance professionals perceive Digital Engineering capabilities. Fourth, compliance frameworks and IT service management standards explicitly require capabilities that Digital Engineering provides. Yet the academic community has not investigated whether Digital Engineering could fulfill these requirements more effectively than the traditional approaches that frequently fail. Throughout, this review takes as given the conceptual definitions of Digital Engineering and its four pillars established in Chapter~1. The review builds on that foundation to weigh evidence, debate, and gaps rather than to restate definitions. Where a pillar reappears here, it does so as the object of evidentiary appraisal.

\section{Parallel Paths: The Independent Evolution of Adjacent Disciplines}

Understanding why that gap opened, and why it has proved so durable, requires understanding how these disciplines evolved along separate trajectories. Each developed its own academic identity, professional communities, curricular standards, and institutional frameworks, even while confronting similar challenges in documentation, traceability, configuration management, and lifecycle governance. This independent evolution created disciplinary silos so deeply embedded that researchers within each tradition rarely encounter the methods of adjacent fields. The result is a structural consequence of how academic disciplines form and maintain boundaries, an outcome of institutional design working as it normally does.

Systems engineering emerged as a recognizable discipline in the mid-twentieth century, rooted in large-scale military and industrial programs that demanded coordination across multiple engineering specialties. The discipline developed across three historical intervals: foundational (1926--1960), exploratory (1961--1989), and contemporary (1990--2017)~\parencite{Hossain_2020}. Across these intervals, systems engineering ``is not affirmed by a set of underlying fundamental propositions'' but rather emerged as ``a set of best practices'' consolidated from defense and aerospace experience~\parencite{Hossain_2020}. Hall's foundational \textit{A Methodology for Systems Engineering} (1962) and Goode and Machol's \textit{System Engineering} (1957) established the intellectual architecture, while organizations such as Bell Telephone Laboratories and the RAND Corporation provided the institutional incubators~\parencite{Hossain_2020}. Reviewing seventy years of the discipline's existence, \citeauthor{Honour_2018} noted that systems engineering experienced ``rapid growth'' in its early decades followed by ``partial stagnation in the 1990s and 2000s'' as prescriptive process standards displaced first-principles engineering thinking~\parencite{Honour_2018}. This stagnation coincided with the period when enterprise IT complexity was accelerating. The temporal mismatch helps explain why IT organizations did not adopt systems engineering methods during the critical window when they might have been most needed.

Disciplinary insularity deepened as systems engineering grew more abstract. A companion history traced the evolution from a technical-objective focus through a process and methodology focus toward contemporary calls to return to foundational roots, arguing that systems engineering ``became more abstract and disconnected from the rest of engineering'' when prescriptive processes were developed~\parencite{Pennotti_2024}. As part of INCOSE's Future of Systems Engineering initiative, that same work characterized systems engineering as inherently transdisciplinary yet insufficiently recognized as such. This characterization carries direct implications for the present research. If systems engineering's own professional community acknowledges that the discipline has been too insular, the failure to extend its methods to enterprise IT becomes explicable. It stems from disciplinary boundaries that neither community has been motivated to cross, boundaries left in place by default while the question of applicability went unasked. Survey evidence corroborates this self-assessment: across INCOSE chapter attendees in North America, a measurable gap separates systems engineering as a documented discipline from its actual technical practice~\parencite{Wasson_2019}.

Crucially, throughout these seven decades, systems engineering oriented itself toward physical systems: aircraft, naval vessels, spacecraft, weapons platforms, and industrial control systems. Enterprise information systems---encompassing servers, networks, databases, cloud infrastructure, and software applications---fell outside the discipline's traditional scope. Yet these systems grew to exhibit the same characteristics of complexity, interdependency, and emergent behavior that justified systems engineering approaches for physical platforms. The Systems Engineering Body of Knowledge acknowledges this historical trajectory, documenting that systems engineering principles apply to ``engineered systems'' broadly while noting that most published case studies and validated methodologies derive from defense and aerospace contexts~\parencite{SEBoK_2024}.

Information Assurance, by contrast, developed from different institutional origins. Where systems engineering emerged from engineering programs and defense laboratories, Information Assurance developed at the intersection of national security policy and computing education. \citeauthor{Maconachy_2001} formalized the discipline's conceptual foundations in their influential model extending McCumber's 1991 information security cube to encompass five security services (confidentiality, integrity, availability, authentication, and non-repudiation) across three information states and three countermeasure categories~\parencite{Maconachy_2001}. The framework emerged under the auspices of the National Security Agency and Idaho State University's National Information Assurance Training and Education Center. It established Information Assurance as a discipline concerned with protecting and governing information across its lifecycle rather than with engineering the systems that process it. The Committee on National Security Systems codified training standards through NSTISSI documents 4011 through 4016, and the Centers of Academic Excellence in Information Assurance Education program, established in 1998, formalized the academic pipeline~\parencite{Dark_2006}. Information Assurance thus became a ``pervasive theme'' within IT curricula, embedded within computing education but distinct from the engineering disciplines where systems engineering resided~\parencite{Dark_2006}.

A comprehensive evaluation of Information Assurance and security models traced how the field's models evolved, moving from McCumber's original cube through the Maconachy-Schou-Ragsdale model and arriving at a new Reference Model of Information Assurance and Security. That evaluation drew on twenty-six expert interviews, three workshops, and a case study~\parencite{Cherdantseva_2016}. That analysis revealed that existing IA models addressed ``trends of diversification and de-perimeterisation'' but focused on security properties and countermeasures rather than on the engineering lifecycle processes through which systems are designed, built, and maintained. This focus on protection properties rather than engineering processes is a fundamental orientation difference from systems engineering, and it explains why Information Assurance professionals have not encountered Model-Based Systems Engineering in their professional formation.

Cybersecurity's emergence as a standalone academic discipline further fragmented the intellectual landscape. \citeauthor{Parrish_2018} characterized cybersecurity as a ``meta-discipline,'' an aggregate label encompassing contributions from computer science, information systems, information technology, software engineering, and electrical engineering~\parencite{Parrish_2018}. Their analysis traced how establishing the CAE/IAE program in 1998 initiated institutional recognition that evolved from Information Assurance education into cybersecurity education over two decades. The Joint Task Force on Cybersecurity Education formalized this disciplinary independence through \textit{Cybersecurity Curricula 2017}, the first global curricular guidelines for post-secondary cybersecurity degree programs, organized around eight knowledge areas with contributions from over three hundred advisors across thirty-five countries~\parencite{CSEC2017}. The committee report acknowledged that ``the field of cybersecurity is in its formative stages''~\parencite{CSEC2017}. Such a characterization positions cybersecurity as a discipline still establishing its boundaries rather than one prepared to absorb methods from engineering fields with which it shares no curricular lineage.

An analysis of master's programs at leading universities examined the evolution of cybersecurity as an academic discipline, finding that it ``is considered to be evolving at a rapid pace'' compared to other scientific and engineering disciplines~\parencite{Cabaj_2018}. The \textit{Computing Curricula 2020} report, the ACM and IEEE Computer Society's comprehensive overview of computing education, positioned cybersecurity alongside established computing disciplines while documenting how each maintained distinct identities despite shared foundations~\parencite{CC2020}. Chapter 2 of that report traced how computing disciplines evolved since CC2005, revealing a pattern of increasing specialization and decreasing cross-disciplinary integration. Systems engineering does not appear among the computing disciplines at all; it occupies a separate academic lineage rooted in industrial and mechanical engineering departments rather than computing schools. This institutional separation means that cybersecurity and IT graduates rarely encounter systems engineering coursework, and systems engineering graduates rarely encounter Information Assurance or IT service management content.

IT Service Management developed along yet another independent trajectory, one driven by practitioner communities rather than academic research. The foundational systematic literature review of ITSM implementation examined research published between 2000 and 2012~\parencite{Iden_2013}. Its findings documented that the earliest academic study addressing ITIL adoption appeared in 2005 and the first journal article in 2006, more than fifteen years after ITIL's initial development by the United Kingdom's Central Computer and Telecommunications Agency in the 1980s. This extended gap between practitioner adoption and academic investigation produced a field with extensive operational guidance but limited theoretical grounding. The most recent comprehensive systematic review of ITSM reinforced this assessment, covering 126 empirical studies from 2012 through 2021 and reconceptualizing ITSM as a Management Control System~\parencite{MacLean_2023}. Their observation that ITSM has been ``undertheorized by academics'' despite widespread practitioner adoption encapsulates the practitioner-academic gap that characterizes the field~\parencite{MacLean_2023}.

The consequences of this practitioner-driven evolution include widespread ``conceptual confusion'' across ITSM models such as ITIL, CobiT, BSM, SLM, and ISO/IEC 20000~\parencite{Winniford_2009}. Multiple frameworks addressing overlapping concerns developed without mutual awareness or coordination, a fragmentation pattern that mirrors the broader disciplinary silos this section documents. In a cross-national empirical study of 623 respondents from the United Kingdom, United States, German-speaking countries, and Australia, \citeauthor{Marrone_2014} found that operational ITIL processes achieved substantially higher adoption rates than strategic processes~\parencite{Marrone_2014}. A systematic literature review by \citeauthor{Serrano_2021} synthesized forty-seven articles mapping the ITSM field's evolution since 2002, tracing how ITSM grew from the service management concept emerging in the 1930s through ITIL's development into a distinct professional domain~\parencite{Serrano_2021}. Throughout this evolution, ITSM developed no connection to systems engineering practice. ITIL's process-oriented framework emerged from operational IT management experience rather than from engineering design principles, creating a fundamental methodological incompatibility: ITSM treats IT services as processes to be managed, while systems engineering treats systems as engineered artifacts to be designed, verified, and maintained through lifecycle governance.

These parallel evolutionary paths produced more than an absence of applied research. The deeper deficit is that no intellectual infrastructure exists to enable such application. Researchers in each discipline publish in different journals, attend different conferences, cite different foundational theorists, and frame problems through different conceptual lenses. A systems engineer describes configuration management as maintaining the consistency of a system's physical and functional attributes with its requirements and design documentation~\parencite{SEBoK_2024}. An IT service manager describes configuration management as maintaining information about configuration items required to deliver IT services~\parencite{ITIL_4_2019}. An Information Assurance professional describes configuration management as establishing and maintaining the integrity of systems through control of changes~\parencite{NIST_SP_800_128}. These three descriptions address the same operational challenge through three different conceptual vocabularies, each embedded within frameworks that do not reference one another. The research gap this dissertation investigates is, at its root, a gap between disciplines that never developed the shared vocabulary necessary for methodological exchange.

\subsection{A Theoretical Lens on the Divide: Boundary Objects and Knowledge Transformation}

Why a divide of this kind should prove so stubborn, and how it might be crossed, comes into sharper focus when the gap is viewed through the theory of knowledge boundaries. The pivotal idea is the \emph{boundary object}, introduced by \citeauthor{Star_Griesemer_1989} in a study of how naturalists, collectors, and administrators cooperated productively without sharing a common expertise. A boundary object is an artifact ``plastic enough to adapt to local needs and the constraints of the several parties employing them, yet robust enough to maintain a common identity across sites''~\parencite{Star_Griesemer_1989}. In plainer terms, it is a shared representation---such as a map, a specification, or a model---around which communities that speak different professional languages can coordinate their work without first agreeing on terms. The concept names precisely the sort of artifact that might let a systems engineer, a security analyst, and an IT service manager reason about the same system together.

Why the present divide resists easy bridging emerges from the refinement that \citeauthor{Carlile_2004} brought to the idea. He graded knowledge boundaries by the difficulty of crossing them: syntactic boundaries yield to the simple \emph{transfer} of information, semantic boundaries demand the \emph{translation} of differing interpretations, and pragmatic boundaries, where the parties hold genuinely different interests and entrenched practices, yield only to the costly \emph{transformation} of knowledge itself~\parencite{Carlile_2002, Carlile_2004}. The boundary separating systems engineering from Information Assurance and IT service management is pragmatic in exactly this sense. The three vocabularies of configuration management cataloged above are not a shortfall of translation that a shared glossary could repair; they encode different practices, incentives, and professional identities, and bridging them therefore demands the transformation of working knowledge rather than the mere exchange of definitions.

Several complementary lenses deepen this account of why the divide endures. \citeauthor{Cohen_Levinthal_1990} established that a community's capacity to absorb external knowledge depends on its stock of prior related knowledge~\parencite{Cohen_Levinthal_1990}. Disciplines that have never encountered one another's methods therefore have little of the absorptive capacity required to assimilate them. \citeauthor{Szulanski_1996} showed empirically that the principal impediments to transferring proven practice are knowledge-related rather than motivational, a phenomenon he termed internal stickiness~\parencite{Szulanski_1996}. \citeauthor{Brown_Duguid_1991} and \citeauthor{Wenger_1998} located knowledge within communities of practice, whose tacit and practice-bound knowing resists movement across community boundaries~\parencite{Brown_Duguid_1991, Wenger_1998}. At a broader scale, \citeauthor{Gibbons_1994} characterized a shift toward knowledge produced in the context of application by transdisciplinary teams~\parencite{Gibbons_1994}, and \citeauthor{Bammer_2020} argue that the integration expertise needed to span such teams is itself a distinct and frequently absent capability~\parencite{Bammer_2020}. Read together, these perspectives recast the systems-engineering-to-information-assurance gap not as a failure of effort but as the predictable behavior of distinct epistemic communities, each holding limited absorptive capacity for the other's practice-bound knowledge.

That such transformation is arduous yet disproportionately valuable is the central message of the National Academies' assessment of interdisciplinary work. The assessment defines the enterprise as integrating ``information, data, techniques, tools, perspectives, concepts, and/or theories from two or more disciplines'' to attack problems ``beyond the scope of a single discipline,'' and it catalogs the institutional barriers that keep such integration rare~\parencite{NRC_Interdisciplinary_2005}. Systems engineering has itself responded to this call. \citeauthor{Madni_TSE_2018} reframes the discipline as inherently transdisciplinary, urging it to reach beyond its traditional borders and to exploit convergence with adjacent fields rather than retreat behind them~\parencite{Madni_TSE_2018, Madni_2019_TSE}. Nor is the maneuver merely aspirational. Working in the security-critical setting of collaborative cyber-physical systems, \citeauthor{Ramli_2021_BoundaryObjects} treat a reference architecture as a boundary object that aligns otherwise divergent stakeholders around shared trustworthiness properties~\parencite{Ramli_2021_BoundaryObjects}. This illustrative application suggests that the concept can perform real work at the frontier between engineering and security.

Read together, these strands recast the research gap as something deeper than an empirical absence. Theory anticipates both the difficulty and the remedy. The divide persists as the predictable behavior of a pragmatic knowledge boundary, and the most promising instrument for crossing such a boundary is precisely the shared, authoritative, machine-readable model that Digital Engineering produces. Whether the models of Digital Engineering can in fact serve as the boundary objects that carry transformed knowledge between systems engineering, Information Assurance, and IT service management is, at bottom, the question this dissertation begins to answer. The present study takes only the first step toward that answer, and a deliberately bounded one. It measures whether the professionals who would have to cross this boundary perceive value in the capabilities that might carry them across it, which corresponds to the perceived-usefulness component of the adoption theories introduced in Chapter~1. Whether Digital Engineering models actually function as boundary objects, transforming knowledge across these communities in practice, is a question a perception survey cannot settle. This study therefore frames the boundary-object account as the theory that the implementation and longitudinal studies proposed in Chapter~5 are designed to test, rather than as a claim the present data establish.

\section{The Evidence Paradox: Demonstrated Value Without Sufficient Proof}

Model-Based Systems Engineering rests on a paradox that shapes its entire adoption landscape: practitioners and organizations widely perceive MBSE as valuable, yet rigorous empirical measurement of that value remains exceptionally rare. Understanding this paradox is essential because it explains both why MBSE has succeeded within systems engineering and why it has not expanded beyond it. Organizations already embedded in the systems engineering discipline can rely on institutional experience and professional networks to validate MBSE investment. Organizations outside the discipline, including enterprise IT and Information Assurance, lack comparable experiential evidence and must rely on published research that provides insufficient quantitative justification.

The most comprehensive assessment of MBSE evidence, a systematic literature review, categorizes reported benefits into four evidence types: measured, observed, perceived, and referenced~\parencite{Henderson_Salado_2021}. Approximately two-thirds of claimed MBSE benefits lack empirical measurement and instead rely on perceived or referenced evidence, a finding with direct consequences for adoption decisions. A systematic review by \citeauthor{Carroll_Malins_2016} reached a convergent conclusion from a separate corpus: of 88 MBSE case studies examined, 67 reported no quantitative metrics, and the authors cautioned that case-study author bias tends to distort the return-on-investment figures that the remainder do report~\parencite{Carroll_Malins_2016}. Organizations evaluating MBSE investment discover that the research literature offers abundant testimony to MBSE's promise but limited quantified outcomes on which to build business cases. A review of MBSE motivation and research opportunities reached a similar conclusion, arguing that the value proposition requires demonstration through real-world applications and that further advancements remain necessary before broader adoption can occur~\parencite{Madni_Sievers_2018}. The more recent systematic literature review by \citeauthor{Wooley_Womack_2025} analyzed adoption, benefits, and challenges across the Digital Engineering research corpus and returned no research addressing enterprise IT infrastructure or Information Assurance applications~\parencite{Wooley_Womack_2025}. This finding corroborates that the gap reflects the state of the literature rather than an incomplete search.

Among the rare exceptions to the measurement deficit, the work of \citeauthor{Rogers_Mitchell_2021} stands as the most frequently cited quantified evidence of MBSE return on investment~\parencite{Rogers_Mitchell_2021}. Their case study examined the Submarine Warfare Federated Tactical Systems program, a rapidly evolving combat system of systems comprising over ten million source lines of code deployed across 104 submarines, following a \$3 million MBSE investment over two and a half years. A like-for-like comparison between legacy document-centric processes and the MBSE approach revealed that systems engineering hours per requirement decreased from 12.1 to 9.9, representing an 18\% improvement in efficiency that exceeded the 13\% improvement projected by an earlier pilot study. The MBSE approach handled 42\% more interface requirements changes while consuming only 16\% more hours, and it reduced total interface defects by 9\%. It also shifted 18\% of defects to earlier discovery phases, where correcting a defect costs roughly one-fourth to two-thirds of what correction costs during platform integration testing. Supplementary data demonstrated 30\% more baselines produced monthly and 60\% more integrated subsystems and combat system variants managed within constant resources~\parencite{Rogers_Mitchell_2021}.

These results stand out for their scarcity as much as for their magnitude. The same evidence assessment identified the Rogers and Mitchell study as one of only two papers in the literature reporting measured MBSE return on investment evidence~\parencite{Henderson_Salado_2021}. The intervening years have added only sparingly to that record, and the most recent addition warrants both caution and attention. In a 2025 mixed-methods study of the post-production design-change process for United States Navy assets, \citeauthor{Shefa_Topcu_2025} estimated that a Digital Engineering transformation could reduce median project duration by roughly half while sharply lowering its variability, although the gains varied across task categories~\parencite{Shefa_Topcu_2025}. Two cautions attach to this finding: it remains a preprint awaiting peer review, and its setting is naval sustainment rather than enterprise IT. Neither caution is disqualifying, and the study earns its place on three grounds. Its mixed-methods design triangulates quantitative schedule data against practitioner interviews rather than resting on testimony alone; the magnitude it reports is consistent with the rare earlier measurement from \citeauthor{Rogers_Mitchell_2021}; and its object of study, a recurring change-and-configuration process burdened by document-centric handoffs, is structurally close to the change management that consumes Information Assurance and IT service management teams. Treated as corroboration rather than as decisive proof, the result reinforces the governing pattern: where measured returns appear at all, they tend to be large, and what stays scarce is the measurement itself.

Subsequent efforts to address this measurement deficit have advanced the methodological conversation without resolving it. A measurement framework for selecting and developing metrics to assess MBSE and Digital Engineering value included causal maps of expected benefits~\parencite{Henderson_2023_Metrics}. A separate evaluation of MBSE's perceived value analyzed attributes, benefits, and drawbacks reported in the literature, finding a persistent gap between what practitioners believe MBSE delivers and what published evidence can substantiate~\parencite{Campo_2023}. Both studies underscore a recurring theme. The systems engineering community recognizes the measurement deficit and has begun developing frameworks to address it, yet empirical data from operational programs remain rare. Authoritative recent assessments continue to refine, rather than resolve, the measurement question. Two RAND industry studies conclude that the benefits of Digital Engineering are genuine but accrue over the long term and resist immediate cost or schedule quantification~\parencite{RAND_DE_Supply_2024, RAND_DE_CostBenefit_2024}, and a Defense Science Board task force similarly reports demonstrated lifecycle benefits set against substantial workforce and tooling costs~\parencite{DSB_DE_2024}. The multi-organization Digital Engineering Measurement Framework responds directly to this difficulty by defining standardized benefit areas against which value can be assessed~\parencite{PSM_DE_Measurement_2022}. Survey evidence reinforces the perception dimension, as \citeauthor{Brown_Zaidi_2025} found organizational resistance and a shortage of trained engineers to be the dominant adoption barriers among aerospace and defense practitioners~\parencite{Brown_Zaidi_2025}.

For enterprise IT organizations operating on annual budgets with continuous delivery expectations, this evidentiary deficit creates a rational basis for hesitation: the demonstrated benefits derive from a naval combat system program operating under a level-of-effort contract structure fundamentally different from enterprise IT budget models. Whether comparable returns would materialize in enterprise IT remains untested.

Earlier systematic reviews reinforce this pattern. A review of thirteen years of SysML research mapped the landscape of model-based methods and their empirical evidence~\parencite{Wolny_2020}. A survey of MBSE tools and applications found consistent reports of improved requirements traceability and stakeholder communication alongside persistent concerns about tool complexity and organizational adoption challenges~\parencite{Chami_Bruel_2018}. Collectively, these reviews establish that MBSE has demonstrated value across multiple domains while simultaneously documenting two critical limitations: the evidence remains predominantly qualitative, and the domains of application remain overwhelmingly confined to aerospace and defense.

The tension between demonstrated capability and insufficient proof creates what might be termed an adoption credibility gap. Within the systems engineering discipline, professional experience and institutional knowledge compensate for the measurement deficit. Systems engineers who have used MBSE tools can observe improvements in their own work, creating the experiential validation that published evidence lacks. Outside the discipline, this compensatory mechanism does not operate. IT and Information Assurance professionals have no experiential basis for evaluating MBSE claims. The literature provides them with compelling arguments but insufficient evidence, a combination that technology adoption research suggests is inadequate for driving adoption decisions.

\section{Adoption as a Perception Problem}

If the evidence paradox explains why MBSE has not expanded through rational economic calculation, the adoption literature reveals a complementary explanation rooted in perception dynamics. A growing body of empirical research shows that MBSE adoption decisions are driven more by how professionals perceive the technology than by its objective characteristics. This finding carries direct implications for the present research: if perceptions govern adoption even among systems engineering professionals who work with models daily, understanding perceptions among IT and Information Assurance professionals, who have never systematically encountered these methods, becomes a prerequisite for any meaningful discussion of cross-disciplinary transfer.

\citeauthor{Call_Herber_2022} supply the theoretical foundation for this argument, mapping Everett Rogers' Diffusion of Innovations theory to MBSE adoption dynamics~\parencite{Call_Herber_2022}. Rogers' framework identifies five perceived attributes of innovations (relative advantage, compatibility, complexity, trialability, and observability) that collectively account for 49--87\% of variance in adoption rates across innovation research. The critical theoretical insight, and the one most consequential for the present research, is that \textit{perceptions} of these attributes, not the attributes themselves, drive adoption behavior. The authors invoke W.~I.~Thomas's dictum that ``if [people] define situations as real, they are real in their consequences,'' arguing that shaping how MBSE attributes are perceived can accelerate adoption rates more effectively than improving the attributes themselves~\parencite{Call_Herber_2022}.

An empirical follow-up by \citeauthor{Call_2024} tested this theoretical mapping through a survey distributed to systems engineering professionals via INCOSE and professional networks, which drew 270 responses, 197 of them complete~\parencite{Call_2024}. The study examined perceptions across six respondent subpopulations, among them those involved versus not involved in MBSE efforts, model users versus non-users, and respondents subject to Digital Engineering mandates versus those who are not. Respondents broadly recognized MBSE's relative advantage in improving data quality and traceability. However, perceived compatibility with existing practices and perceived complexity emerged as significant barriers that suppressed adoption even when practitioners acknowledged objective value.

Most striking is the finding on trialability. Respondents not involved in MBSE efforts reported far lower access to trial opportunities and tools; for the statement ``I have opportunities to experiment with MBSE,'' the chi-squared test of independence between the involved and non-involved subpopulations produced the largest divergence in the study, with \(\chi^2 = 72.1\) and reported \(p = .000\)~\parencite{Call_2024, Call_2024_Dissertation}.
% Verified 2026-07 against Call's CSU dissertation (author version, engr.colostate.edu/~drherber/files/Call2024b.pdf): survey received 270 responses, 197 complete (33 blank), open Jun--Sep 2023. Six subpopulations analyzed (Appendix C: Involvement with MBSE, Model Use, Organization MBSE Use, Organization SE Plan, Subject to DE Mandates, Associated with Air Force); the three named here are examples. Q14 involved/non-involved comparison reported as chi-squared with p = 0.000 (Table 5.1), chi-squared = 72.1 (Appendix C, Table C.1); symposium paper reports p = 0.00. Degrees of freedom not reported in either source. "Preventative innovation" framing confirmed (Call ch. on relative advantage). Bib keys Call_2024 (pages 462--478) and Call_2024_Dissertation both present and correct.
This creates what \citeauthor{Call_2024} characterize as a barrier-reinforcing cycle: professionals who have not experienced MBSE cannot access the trial opportunities that would enable informed adoption decisions, while those who have experienced MBSE report substantially more favorable perceptions. The study further characterizes MBSE as a ``preventative innovation'' whose advantages derive from preventing problems such as inconsistencies, documentation errors, and rework rather than from producing visible new benefits. This characteristic depresses both perceived relative advantage and observability, compounding adoption barriers even when objective value exists.

This perception-adoption dynamic has profound implications for cross-disciplinary transfer. If systems engineering professionals, practitioners who work within the discipline that developed MBSE, exhibit perception gaps that impede adoption, the barriers among IT and Information Assurance professionals are likely to be greater still on the attributes that depend on contact with the methods, namely trialability and observability. Relative advantage may behave differently. Digital Engineering addresses documentation, traceability, and visibility problems these professionals already face, so a favorable reading of its advantage could survive transfer even where exposure is absent, which is the expectation Proposition~P2 states and Chapter~4 tests. These professionals operate outside the systems engineering discipline entirely. They have no professional exposure to MBSE concepts, no institutional networks sharing MBSE experiences, and no trial opportunities through which favorable perceptions might develop. The present research addresses this gap by measuring perceptions in a population where no empirical data currently exist.

Complementary evidence from \citeauthor{Henderson_McDermott_Salado_2024} deepens understanding of adoption dynamics through a mixed-methods study combining systematic literature review with eighteen semi-structured practitioner interviews from organizations attempting MBSE adoption~\parencite{Henderson_McDermott_Salado_2024}. The literature review extracted over 600 individual lessons learned from 46 published papers, coded into categories spanning communication, model definition, organizational strategy, and technical implementation. The interview findings illuminate the human dimensions of adoption that quantitative studies cannot fully capture. Most directly relevant to the question of professional awareness, approximately 22\% of interview participants, all of them systems engineering professionals, could not convey a clear definition of MBSE\@. One participant captured the prevailing confusion: ``Most people think of MBSE as being synonymous with a specific tool'' and ``don't understand how MBSE fits into DE or what it really does other than it is modeling instead of documents''~\parencite{Henderson_McDermott_Salado_2024}. Another stated plainly that ``if they aren't familiar with MBSE, they're not going to use it.''

Organizational lessons proved equally revealing. For both organizations in the study whose adoption efforts failed, the driving barrier was lack of management support. Middle management resistance proved especially damaging because executives endorsed MBSE while middle managers could disregard the endorsement without operational consequences. Successful adopters commonly established core MBSE teams, such as communities of practice or centers of excellence, and implemented role-based training at four levels: model reviewers for leaders and decision-makers, developers and modelers, architects for senior engineers, and administrators for IT and software support. A reinforcing dynamic emerged: greater stakeholder exposure to models produced more perceived benefits, which generated more organizational buy-in, which enabled further exposure. This virtuous cycle operates within organizations that have initiated adoption but cannot initiate itself in domains where MBSE has never been introduced. A synthesis of five years of research on Digital Engineering and MBSE adoption produced an enterprise adoption framework, a model of factors spanning organizational design, enablers and barriers, and change management~\parencite{McDermott_2024}. The synthesis confirmed that adoption requires coordinated action across multiple organizational dimensions simultaneously, a finding that compounds the challenge for enterprise IT organizations unfamiliar with the discipline.

A survey of 51 practitioners extended these findings, examining how organizational structure variables correlate with MBSE adoption outcomes~\parencite{Henderson_Salado_2024_Org}. The results reveal a pattern with direct implications for enterprise IT contexts. Flexibility and interconnectedness showed the strongest and most consistent positive correlations with both adoption process and implementation variables. Formalization, defined as documented procedures and processes, correlated positively, a somewhat counterintuitive finding suggesting that governance frameworks support rather than hinder adoption when combined with flexibility. Centralization correlated negatively with adoption and implementation, as did large organizational size and high vertical differentiation. These structural findings suggest that enterprise IT organizations characterized by rigid hierarchies, centralized governance, and siloed operations may face structural impediments to Digital Engineering adoption that compound the perception barriers documented by \citeauthor{Call_2024}.

Research beyond the systems engineering discipline corroborates and extends these dynamics. The embedded systems industry sits outside the defense and aerospace context where MBSE originated. A qualitative study of twenty interviews across ten companies in that sector identified immature tooling, return on investment uncertainty, and migration fears as key barriers, confirming that the adoption challenges documented within systems engineering intensify when the technology crosses disciplinary boundaries~\parencite{Vogelsang_2017}. A broader study of strategic adoption of digital innovations found that digital transformation requires coordinated enterprise-level application rather than bottom-up adoption of individual technologies~\parencite{Campagna_2024}. The research identified twelve strategic adoption influencers and noted that adoption research focuses on individual technologies rather than integrated digital transformation. This finding explains why platform-level Digital Engineering adoption within defense programs has not expanded to enterprise IT functions operating separately from program organizations.

Technology adoption research in cybersecurity and IT contexts provides a parallel body of evidence that has not been connected to the Digital Engineering adoption discourse. One study integrated Diffusion of Innovation, the Technology Acceptance Model, and the Technology-Organization-Environment framework with a balanced scorecard approach, surveying 147 IT experts in United Kingdom small and medium enterprises to identify factors affecting cybersecurity technology adoption~\parencite{Hasani_2023}. Their study is one of the few empirical investigations combining multiple adoption models in a cybersecurity context, identifying eight key factors including organizational IT readiness, perceived ease of use, and management support. These factors parallel the MBSE adoption barriers documented by \citeauthor{Henderson_McDermott_Salado_2024}. In a related effort, \citeauthor{Alghamdi_2023} developed a cybersecurity-specific extension of the Unified Theory of Acceptance and Use of Technology, surveying 108 IT professionals and 554 public respondents to examine cybersecurity adoption in smart city contexts~\parencite{Alghamdi_2023}. Drawing on the same theoretical foundation that \citeauthor{Call_2024} used for MBSE, \citeauthor{Anthony_2024} applied Rogers' Diffusion of Innovations framework to examine motivations and barriers for advanced cybersecurity tool adoption in high-threat industries~\parencite{Anthony_2024}. That cybersecurity researchers independently selected the same adoption theory applied by MBSE adoption researchers underscores the theoretical compatibility between these fields despite their practical isolation.

At the organizational level, a large-scale empirical study using data from 1,988 company executives examined IT innovation adoption factors, finding that organizational IT innovation readiness, which encompasses both technology readiness and internal expertise, is a key mediator of adoption outcomes~\parencite{Ali_2022}. A study combining Diffusion of Innovation and Technology-Organization-Environment frameworks surveyed 535 managers and found that relative advantage, complexity, observability, top management support, and competitive pressure serve as significant determinants of digital technology adoption~\parencite{Shahadat_2023}. These organizational-level findings parallel the structural variables examined by \citeauthor{Henderson_Salado_2024_Org} in the MBSE context, suggesting that common adoption dynamics operate across technology domains even when the specific technologies differ.

The adoption literature thus establishes a critical precondition for the present research. Perceptions govern technology adoption. Even among systems engineering professionals, the population most favorably positioned to appreciate MBSE, perceptions remain mixed and marked by substantial awareness deficits. Cybersecurity and IT adoption research reveals parallel dynamics operating through the same theoretical frameworks. Yet no study bridges these parallel tracks by examining Digital Engineering perceptions among IT and Information Assurance professionals. Until such data are collected, any discussion of Digital Engineering adoption beyond defense and aerospace rests on assumption rather than evidence. This dissertation addresses that evidentiary gap.

\section{Crossing Disciplinary Boundaries: Digital Engineering Beyond Its Origins}

Digital Engineering's concentration within aerospace and defense raises an unresolved question that the literature has acknowledged but not adequately investigated: can these methods transfer effectively to domains with different professional cultures, tool ecosystems, and economic structures? The preliminary evidence is suggestive but fragmented, consisting of isolated applications to security domains, a conceptual framework from an authoritative institution, emerging applications in healthcare and manufacturing, and a single systematic mapping study of model-based security engineering. No sustained research program has examined cross-disciplinary transfer to enterprise IT or Information Assurance.

\citeauthor{Bone_Blackburn_2019} described the Systems Engineering Research Center's research effort underlying the Department of Defense Digital Engineering initiative, identifying three key enablers for transformation: IT infrastructure, workforce development, and policy~\parencite{Bone_Blackburn_2019}. Though defense-originated, their analysis addressed cross-cutting adoption dynamics and workforce transformation challenges applicable to any domain adopting Digital Engineering. The transformation they described requires coordinated investment across technology, people, and governance, precisely the integrated approach that enterprise IT organizations, structured around operational silos and annual budget cycles, struggle to sustain.

Among existing efforts, the most substantive bridge between Digital Engineering and Information Assurance comes from the Carnegie Mellon Software Engineering Institute. \citeauthor{CMU_SEI_DevSecOps_2023} argued that DevSecOps pipelines constitute complex socio-technical systems that require systems engineering treatment, comprising components that are independently developed and maintained, physically and logically distributed, and interoperable~\parencite{CMU_SEI_DevSecOps_2023}. The tight integration of business mission, capability delivery, and products ``increases the attack surface of the product under development.'' As organizations adopt DevSecOps tools and techniques with increased coupling between the product being built and the tools used to build it, ``the attack surface of the product continues to grow, incorporating segments of the development environment.'' The research developed a DevSecOps Platform Independent Model using the Unified Architecture Framework and SysML to model pipeline security properties, demonstrating how to construct cybersecurity assurance cases from model-based representations.

What gives such a connection its analytical rigor is the assurance case, a connective concept with an established lineage of its own. An assurance case is a structured, evidence-supported argument that a system satisfies a claimed property such as safety or security. \citeauthor{Kelly_Weaver_2004} formalized its dominant notation, the Goal Structuring Notation, which renders the argument as an explicit hierarchy of claims, supporting strategies, and the evidence that discharges them~\parencite{Kelly_Weaver_2004}. \citeauthor{Wei_2019_SACM} subsequently carried the notation into the model-based world through the Object Management Group's Structured Assurance Case Metamodel, enabling assurance arguments to be generated and maintained as formal models rather than as static documents~\parencite{Wei_2019_SACM}. The consequence for Information Assurance is immediate and concrete. When the relationships among requirements, controls, and evidence already reside inside a Digital Engineering model, the compliance argument that auditors demand can be assembled from that model rather than reconstructed by hand. That manual reconstruction is the precise labor that document-centric practice imposes and that the compliance frameworks examined later in this chapter render ever more onerous.

Most relevant to this dissertation, the SEI framework proposed a conceptual shift from process-based to property-based security assurance. Traditional cybersecurity assurance relies on process-based standards such as the NIST Risk Management Framework and SP~800-53 security controls. As systems become more complicated and interconnected, \citeauthor{CMU_SEI_DevSecOps_2023} argued, ``process-based standards fail to assure system owners that the system functions only as intended under all operational circumstances''~\parencite{CMU_SEI_DevSecOps_2023}. When model-informed threat modeling and security practices are done well, ``the overall risks associated with the DevSecOps pipeline and associated products will be reduced, and the compliance and legal requirements will naturally be addressed within the engineering lifecycle.''
% Verified: all four quoted passages confirmed verbatim against CMU/SEI-2023-TR-001 PDF (resources.sei.cmu.edu, DOI 10.1184/R1/22592884), pp. 2-4 and the threat-modeling/assurance sections. The "attack surface ... continues to grow" quote and the "when done well" lead-in were corrected to match the source.
This finding directly supports the theoretical premise of this dissertation: that Digital Engineering capabilities can integrate security engineering into development and operational processes rather than treating compliance as a separate activity conducted after the fact.

However, the SEI framework is a conceptual and architectural argument rather than empirical validation. The research demonstrates technical feasibility through detailed modeling of configuration management capabilities but does not provide quantitative effectiveness data. The framework was designed for heavily regulated and cybersecurity-constrained environments including defense, banking, and healthcare. These are domains where Information Assurance professionals operate and where compliance verification demands consume substantial organizational resources. Lacking empirical validation, the framework leaves the critical question unanswered: does MBSE-based security assurance deliver measurable improvements over traditional compliance approaches in practice?

Beyond the SEI framework, emerging research has begun applying MBSE to non-defense domains, though these applications remain isolated and preliminary. One MBSE methodology for capturing cybersecurity requirements in space systems used SysML within Cameo Systems Modeler, expanding on NIST cyber resiliency frameworks to show how model-based approaches can structure security analysis within an engineering lifecycle~\parencite{Gordon_Reilly_2023}. In a scoping review of MBSE methods applied to rural healthcare system disaster planning, \citeauthor{Berg_2023} demonstrated that model-based approaches can organize complex interdependencies beyond traditional engineering contexts~\parencite{Berg_2023}. A rapid review by \citeauthor{Xames_2024}, covering twelve years of MBSE use in healthcare systems, identified applications ranging from patient flow modeling to wellness digital twins~\parencite{Xames_2024}. A practical MBSE methodology for digital transformation of existing industrial enterprises, developed by \citeauthor{Badenko_2024}, demonstrated feasibility in a non-defense application outside the discipline's traditional domains~\parencite{Badenko_2024}. Reviewing cross-industry principles for digital representations of complex technical systems using the MBSE approach, \citeauthor{Bolshakov_2023} explicitly addressed applications beyond defense and aerospace~\parencite{Bolshakov_2023}. A common pattern runs across these applications: each demonstrates feasibility within a single bounded domain, and none attempts the enterprise-wide integration across infrastructure, security, and service management that this dissertation examines.

Other researchers have probed adjacent territory without directly addressing enterprise IT\@. One MBSE methodology for critical infrastructure vulnerability assessment and decision analysis used DoDAF-based modeling to link regulatory requirements, system architecture, and attack vectors~\parencite{Huff_Medal_2019}. That work demonstrates MBSE application to infrastructure protection and security assessment beyond traditional platform engineering but does not extend to enterprise IT infrastructure management. A UML-based MBSE security profile conforming to ISO/IEC 27001 found, through a feasibility survey of ten engineering companies, that standard SysML and popular MBSE methods inadequately address security aspects~\parencite{Mazeika_Butleris_2020}. This finding identifies a tool gap that would impede adoption even if organizational and perception barriers were overcome: the modeling languages themselves require extension to accommodate the security constructs that enterprise IT and Information Assurance demand. A further proposal, SysML-SecA, combined SysML with security analysis techniques for integrated threat modeling~\parencite{Apvrille_Roudier_2015}. These preliminary investigations establish that researchers have recognized the potential of MBSE for security applications while simultaneously documenting the technical, methodological, and empirical gaps that separate potential from realization.

\citeauthor{Nguyen_Ali_Yue_2017} conducted a systematic mapping of forty-eight primary studies on model-based security engineering for cyber-physical systems~\parencite{Nguyen_Ali_Yue_2017}. The study found that most research uses domain-specific languages or UML, focuses on early lifecycle stages, and addresses threats, attacks, and vulnerabilities generically. This mapping reveals a field in its formative stages, one where foundational concepts are being established but where mature, validated methods for operational security engineering have not yet emerged. The gap between model-based security engineering research and operational Information Assurance practice remains substantial.

In the years since that mapping, model-based security engineering has expanded substantially, though still within engineered platform and cyber-physical contexts. \citeauthor{Mitola_Prys_2024} extend the Digital Engineering framework into a repeatable cyber-hardening process, applying functional mission analysis and zero-trust principles to harden systems against advanced persistent threats~\parencite{Mitola_Prys_2024}. \citeauthor{Gampel_Eveleigh_2025} employ Model-Based Systems Engineering to streamline NIST Risk Management Framework control implementation and the authorization-to-operate decision for industrial control systems~\parencite{Gampel_Eveleigh_2025}, while \citeauthor{Larsen_2024} and \citeauthor{Nolte_2025} pull threat modeling and identifying security requirements forward into the early model-based design of cyber-physical systems~\parencite{Larsen_2024, Nolte_2025}. Extending the assurance-case lineage traced earlier, \citeauthor{Mohamad_CASCADE_2023} present an asset-driven method for constructing security assurance cases aligned to industry standards~\parencite{Mohamad_CASCADE_2023}. These contributions confirm that model-based security engineering continues to advance actively. Yet each remains anchored in engineered platforms, spacecraft, vehicles, and industrial control systems, rather than in enterprise IT infrastructure and Information Assurance programs, which leaves the specific gap this dissertation addresses intact.

Academic literature on MBSE applications thus presents a landscape of expanding interest constrained by limited evidence. Researchers have begun exploring applications to security domains, infrastructure protection, healthcare systems, and compliance frameworks. Authoritative institutions have produced conceptual frameworks demonstrating technical feasibility. Yet the enterprise IT context remains absent from the research, even though organizations in that context manage heterogeneous infrastructure, maintain compliance across multiple regulatory frameworks, and coordinate service delivery across organizational silos. The frameworks exist; the methods have matured; the tools have proliferated. The cybersecurity skills gap analysis by \citeauthor{Ramezan_2023}, examining 935 cybersecurity positions across nine sub-fields, underscores that cybersecurity professionals operate within specialized roles that share no curricular or experiential overlap with systems engineering practice~\parencite{Ramezan_2023}. The research community has not bridged this gap, and the professional communities have not demanded it.

\section{Digital Twins as Emerging Security Tools}

Digital twin technology has generated substantial research interest in cybersecurity applications, creating a parallel track to MBSE-based security engineering that the literature has not yet integrated into a coherent Digital Engineering narrative. The growing body of research on digital twins for security purposes is both an opportunity and an analytical challenge. It demonstrates that Digital Engineering concepts are penetrating security domains, but it does so through a technology-specific lens. That lens misses the integrated approach, combining MBSE, digital threads, digital twins, and PLM, that characterizes Digital Engineering as a discipline.

\citeauthor{Grieves_2023} traces the digital twin concept from its Product Lifecycle Management origins through contemporary applications, positioning the digital twin as integrating physical and virtual systems to enable analysis, optimization, and prediction~\parencite{Grieves_2023}. The foundational distinction between digital twins and traditional simulation lies in synchronization: digital twins maintain continuous data exchange with their physical or logical counterparts, enabling operational decision-making based on current rather than documented system states. An influential early framework for leveraging digital twins in systems engineering established the conceptual architecture that subsequent security applications adapted~\parencite{madni2018leveraging}.

Digital twin security research has matured enough to sustain systematic reviews and comprehensive surveys in recent years. The seminal survey classified security threats across five digital twin components and established the analytical framework that subsequent research has adopted~\parencite{Alcaraz_Lopez_2022}. Its component-level threat taxonomy and systematic coverage of the field make it the reference assessment of the digital twin security landscape. The systematic literature review by \citeauthor{ElHajj_2024} analyzed sixty-seven papers published between 2018 and 2023 examining how digital twins enhance security in Industry 4.0 applications~\parencite{ElHajj_2024}. Covering intrusion detection, vulnerability assessment, cyber range simulation, and threat intelligence applications, the review identified enabling technologies, including machine learning, blockchain, and 5G, used in conjunction with digital twins for security purposes. A comprehensive review extended this analysis, categorizing digital twin cybersecurity studies by technique type and by digital twin level (component, process, asset, system, and network-of-systems) and assessing risk levels from medium to very high depending on twin type and industry sector~\parencite{Alhumam_2025}. A three-dimensional security framework then mapped security risks across digital twin architectural layers and application domains, providing a structured taxonomy for understanding where security challenges concentrate~\parencite{Jeremiah_2024}. A further survey of digital twin modeling techniques classified security operation modes into analytics, simulation, and replication, and reviewed both open-source tools and commercial platforms~\parencite{Qureshi_2025}. These reviews differ in taxonomy and scope, yet they converge on two conclusions: digital twin security applications have achieved sufficient research volume to warrant meta-analysis, and the field's attention remains concentrated on industrial and cyber-physical settings rather than enterprise IT.

Within this expanding literature, several streams bear directly on Information Assurance and enterprise IT applications, and together they carry digital twins across the security lifecycle from offline testing to operational risk assessment. Virtual replicas can support security testing and analysis without affecting operational systems~\parencite{Eckhart_Ekelhart_2019}, train Security Operations Center analysts through a digital twin-based cyber range~\parencite{Vielberth_2021}, and, at the organizational level, help enterprises understand and manage their security postures~\parencite{Dietz_Pernul_2020}. The synchronization that gives twins this reach cuts both ways. Twins can enhance security while simultaneously introducing new attack surfaces~\parencite{Karaarslan_Babiker_2021}, and continuous synchronization between virtual and physical representations enables dynamic threat modeling across the cyber-physical system lifecycle~\parencite{Erceylan_2025}. Progress nonetheless remains bounded. A systematic review of twenty-seven articles on digital twins for incident detection and response in critical infrastructure identified persistent gaps in scalability, interoperability with legacy systems, and real-world validation~\parencite{Kampourakis_2025}. These gaps, particularly interoperability with legacy systems and validation in operational environments, represent precisely the challenges that enterprise IT organizations would confront in digital twin adoption.

Standards development further indicates maturing technology readiness. ISO 23247 and IEC 62832 provide standards for digital twin frameworks~\parencite{Shao_2021}, and subsequent analysis detailed ISO 23247's four-part structure~\parencite{Shao_ISO_23247_2023}. The Internet Engineering Task Force has published a draft reference architecture for network infrastructure digital twins~\parencite{IETF_NDT_2024}. These standardization efforts establish interoperability requirements that reduce vendor lock-in concerns and enable integration across implementations.

NIST's engagement with digital twin technology illuminates the current state of institutional thinking. NIST Internal Report 8356 addresses cybersecurity challenges and trust considerations for digital twin implementations~\parencite{NIST_IR_8356_2025}. Notably, this publication addresses security considerations \textit{for} systems employing digital twins rather than digital twin applications \textit{to} Information Assurance. The distinction is revealing: NIST examines how to secure digital twin implementations rather than how digital twins might enhance security posture visibility or compliance verification. This framing reflects the current state of institutional research: digital twins are treated as systems to be secured rather than as tools for improving security operations.

\subsection{Security Risks of Improperly Secured Digital Twins}

The security of the digital twin itself is directly germane to the enterprise Information Assurance focus of this dissertation rather than a detour into engineering security. A digital twin is itself a critical information repository: an authoritative, continuously synchronized representation of the very systems Information Assurance exists to protect. Like any such repository (the configuration management database is the familiar enterprise example), it must be secured for confidentiality, integrity, and availability. A twin that mirrors production faithfully is a high-value asset squarely within the Information Assurance remit rather than adjacent to it. Its protection is therefore part of the discipline this dissertation addresses, which is why the risks below matter to the argument.

While the preceding discussion examines digital twins as security-enhancing tools, the literature reveals an equally important and less frequently addressed concern: improperly secured digital twins themselves become attack vectors that adversaries can exploit. Because a digital twin faithfully replicates the architecture, configuration, and behavioral characteristics of its physical or logical counterpart, unauthorized access grants an attacker a comprehensive reconnaissance platform. That platform reveals system vulnerabilities, operational patterns, and interdependencies without triggering alerts on the production environment. Framing this danger through the concept of the ``evil digital twin,'' \citeauthor{Suhail_Iqbal_Jurdak_2024} warned that adversaries who compromise or construct malicious virtual replicas can use them to rehearse attacks, identify exploitable weaknesses, and craft intrusion strategies with a precision that traditional reconnaissance cannot match~\parencite{Suhail_Iqbal_Jurdak_2024}. Their analysis identified data poisoning, model manipulation, and feedback loop corruption as three vectors through which compromised twins propagate harm back to physical systems.

Empirical threat analysis corroborates these concerns. \citeauthor{Alcaraz_Lopez_2022}, in the most widely cited survey of digital twin security threats, classified attacks across five architectural components: physical entity, communication channel, data storage, virtual model, and application layer. Each component introduces distinct exploitation opportunities~\parencite{Alcaraz_Lopez_2022}. The resulting taxonomy established that synchronization mechanisms, while essential for twin fidelity, simultaneously create bidirectional pathways through which adversaries can inject false data into the physical system or exfiltrate sensitive operational intelligence from the virtual replica. \citeauthor{Jeremiah_2024} extended this analysis through a three-dimensional security framework mapping risks across digital twin architectural layers and application domains, finding that risk severity escalates from medium to very high as twin fidelity and system-of-systems integration increase~\parencite{Jeremiah_2024}. Higher-fidelity twins, precisely because they model systems more accurately, provide attackers with more actionable intelligence when compromised.

\citeauthor{Holmes_2021} posed the question directly: are digital twins a cybersecurity solution or a cybersecurity challenge? Their analysis concluded that the answer depends on implementation rigor~\parencite{Holmes_2021}. Organizations that deploy digital twins without embedding security controls throughout the twin lifecycle, from design through data ingestion to model synchronization, create shadow representations of critical infrastructure that bypass the security controls protecting the physical systems themselves. \citeauthor{Karaarslan_Babiker_2021} cataloged specific countermeasures including authentication protocols for twin-physical communication channels, encryption of twin state data, and integrity verification mechanisms for synchronization updates~\parencite{Karaarslan_Babiker_2021}. Yet their review found that many implementations prioritize functional fidelity over security hardening, leaving twins exposed to manipulation.

\citeauthor{Suhail_ENIGMA_2023} proposed ENIGMA, an explainable digital twin security solution for cyber-physical systems, specifically addressing the challenge of detecting when a digital twin has been subverted~\parencite{Suhail_ENIGMA_2023}. Their framework employs explainable artificial intelligence to identify anomalous divergences between twin predictions and physical system behavior, divergences that may indicate adversarial manipulation of the twin model. \citeauthor{Jiang_2024_CPPS} demonstrated complementary capabilities in a cyber-physical production systems context, developing a digital twin security model that enhances asset visibility and supports vulnerability prioritization through virtual testing, while simultaneously documenting forty-one vulnerability instances across fourteen critical components that twin-based analysis revealed~\parencite{Jiang_2024_CPPS}.

For Information Assurance practice, these exploitation risks carry particular weight. Organizations contemplating digital twins for compliance verification or security posture management must recognize that the twin itself becomes a high-value target requiring its own security controls, authorization boundaries, and continuous monitoring. These are precisely the governance mechanisms that the RMF prescribes for any information system processing sensitive data. NIST Internal Report 8356 acknowledges this recursive security challenge, establishing trust considerations for digital twin deployments that include data integrity verification, access control requirements, and lifecycle security management~\parencite{NIST_IR_8356_2025}. An unresolved tension persists in the literature: researchers advocate digital twins as tools for improving security visibility while simultaneously documenting that inadequately secured twins expand the very attack surface they are intended to reduce. Resolving this tension requires integrating digital twin security within the broader Digital Engineering framework rather than treating twin security as an isolated concern. Such integration connects twin security to digital thread traceability, MBSE-based security architecture, and PLM lifecycle governance. This integrated approach has yet to be examined in an enterprise IT setting.

Academic research on open source digital twin frameworks adds another dimension to the adoption question. \citeauthor{Gil_2024} conducted a systematic survey of open source digital twin frameworks, analyzing fourteen frameworks against criteria derived from ISO 23247 standards and finding substantial variation in maturity, documentation quality, and community support~\parencite{Gil_2024}. \citeauthor{Autiosalo_Siegel_Tammi_2021} introduced Twinbase, an open source server for the Digital Twin Web concept~\parencite{Autiosalo_Siegel_Tammi_2021}. These open source options reduce economic barriers to adoption but do not resolve the more fundamental challenge: digital twin applications for enterprise IT contexts lack the academic research and validated methods that would guide adoption decisions.

Digital twin security research thus forms a field developing in parallel with but largely disconnected from the broader Digital Engineering discourse. Researchers explore digital twins as security tools without embedding their work within the Digital Engineering framework that integrates MBSE, digital threads, and PLM into a coherent discipline. This disconnection means that researchers investigate individual digital twin security capabilities in isolation rather than as components of the integrated approach that defense and aerospace organizations employ. Whether integrating digital twin security capabilities within a comprehensive Digital Engineering framework would produce greater value than isolated implementations remains an open question, one that the absence of enterprise IT research leaves unaddressed.

\section{The Compliance Imperative and Its Unfulfilled Requirements}\label{sec:compliance_imperative}

Compliance frameworks create requirements that Digital Engineering could address, yet no research identified in this review examines Digital Engineering approaches, as an integrated enterprise discipline, to satisfying these requirements. This disconnect between what frameworks demand and what current practices deliver is one of the most compelling arguments for investigating Digital Engineering application to Information Assurance.

The NIST Risk Management Framework, documented in Special Publication 800-37 Revision 2, provides the authoritative approach to managing security and privacy risk for federal information systems through seven iterative steps: prepare, categorize, select, implement, assess, authorize, and monitor~\parencite{NIST_SP_800_37_R2_2018}. The RMF explicitly requires enterprise architecture integration during the prepare step. Yet compliance with this requirement assumes capabilities that organizations frequently lack: the ability to maintain accurate, current documentation of enterprise architecture that reflects operational reality. NIST Special Publication 800-53 Revision 5 compounds these demands through specific controls requiring enterprise architecture capabilities: PL-2 requires security plans consistent with enterprise architecture; PL-8 requires security architecture development; PM-7 establishes enterprise architecture requirements; CM-2 requires documented baselines; CM-8 requires accurate system component inventory; SA-17 requires design specifications consistent with enterprise architecture~\parencite{Force_2020}. Each control names, in the language of compliance, a capability that Digital Engineering already supplies in its native domain. The mapping between the two vocabularies is what remains undrawn.

The Committee on National Security Systems Instruction 1253 extends these requirements to national security systems, where documentation and visibility challenges compound under additional constraints~\parencite{CNSSI_1253_2022}. Organizations operating outside federal requirements may employ ISO/IEC 27001:2022 for information security management~\parencite{ISO27001-2023}. \citeauthor{Mazeika_Butleris_2020} found that standard MBSE methods inadequately address ISO 27001 security requirements, identifying a specific gap between what compliance frameworks demand and what current modeling approaches provide~\parencite{Mazeika_Butleris_2020}. These alternative frameworks share a common characteristic with NIST guidance: they assume documentation accuracy and visibility capabilities that organizations struggle to maintain.

NIST's systems security engineering publications establish principles that align with Digital Engineering's integrated approach. Special Publication 800-160 Volume 1 Revision 1 describes a basis for engineering trustworthy secure systems~\parencite{Ross_Winstead_McEvilley_2022}. Special Publication 800-160 Volume 2 Revision 1 addresses cyber resiliency considerations~\parencite{NIST_SP_800_160_V2_2021}. These publications reference model-based approaches and traceability requirements without providing specific implementation guidance for Digital Engineering in enterprise IT contexts. The publications establish requirements that Digital Engineering could address: traceability between security requirements, design decisions, and implementation artifacts; documentation that maintains currency throughout system lifecycles; and visibility into system configurations and relationships. Even so, they do not explicitly connect these requirements to Digital Engineering solutions.

Emerging research on compliance automation demonstrates that the academic community has recognized the inadequacy of manual compliance approaches, even if this recognition has not yet connected to the Digital Engineering discourse. Over the past decade this work has coalesced into a coherent research area that treats automated, continuous verification as both feasible and necessary. \citeauthor{Santilli_2023} defined continuous compliance and identified thirteen requirements for automated, ongoing verification from an industrial perspective~\parencite{Santilli_2023}, laying intellectual foundations for a field that did not exist as a coherent research area a decade ago. \citeauthor{Angermeir_2024} extended that foundation with a design science methodology for continuous security compliance in highly regulated domains, showing that automated compliance is both technically feasible and organizationally necessary as system complexity outpaces manual verification capacity~\parencite{Angermeir_2024}. \citeauthor{Castellanos_Ardila_2022} then mapped the landscape through a systematic literature review of compliance checking approaches for software processes, cataloging automated verification methods and identifying research directions~\parencite{Castellanos_Ardila_2022}.

A semantically rich knowledge graph captured regulations from GDPR, PCI DSS, ISO 27001, NIST 800-53, and CSA CCM, enabling automated compliance checking for cloud services~\parencite{Joshi_Elluri_2020}. Validated against the privacy policies of major cloud providers, that work shows that model-based approaches to compliance are technically feasible and practically valuable. The Cloud Property Graph bridged static code analysis and runtime security assessment, using an ontology of cloud resources to identify misconfigurations and regulatory non-compliance across multi-vendor cloud deployments~\parencite{Banse_2021}. A UML-based model for cross-standard security requirements mapped explicitly to NIST's Open Security Controls Assessment Language (OSCAL) Catalog Model, letting organizations track compliance across multiple frameworks simultaneously~\parencite{Stojkov_2021}. Most recently, attack graphs have been combined with OSCAL under compliance-as-code principles to automate risk assessment~\parencite{Koufos_2025}. These efforts differ in formalism, spanning knowledge graphs, property graphs, and UML models, yet they share one architectural commitment: encoding regulatory requirements in machine-readable form so that compliance verification becomes an automated computation rather than a manual clerical task.

The DevSecOps literature provides additional context for this compliance discussion. A systematic review of fifty-four peer-reviewed studies identified twenty-one challenges and thirty-one solutions in DevSecOps adoption, classified into People, Practices, Tools, and Infrastructure themes~\parencite{Rajapakse_2022}. A related analysis identified eighteen challenges mapped to ten core categories, finding that standards exert the most decisive influence on adoption outcomes~\parencite{Akbar_2022}. A multi-vocal literature review identified the primary dimensions of DevSecOps, providing a comprehensive mapping of the field's scope and boundaries~\parencite{Zhao_2024}. Earlier, \citeauthor{Myrbakken_2017} provided one of the first systematic examinations of DevSecOps, defining the field and characterizing the challenges of integrating security into DevOps practices~\parencite{Myrbakken_2017}. These DevSecOps studies establish the problem space of integrating security into development and operational pipelines without connecting to the Digital Engineering framework that could provide structured methodology for that integration.

NIST's own OSCAL initiative represents institutional recognition that manual documentation approaches cannot sustain the accuracy and currency that compliance requires~\parencite{NIST_OSCAL_2023}. OSCAL provides standardized machine-readable formats for expressing security control catalogs, baselines, profiles, and assessment results. The initiative aligns with Digital Engineering's emphasis on machine-readable documentation, and digital thread traceability could connect OSCAL compliance documentation to underlying system configurations, enabling automated verification that documented controls exist as implemented. Yet no academic study has investigated this integration. A notable finding from the literature search is that OSCAL has achieved strong government adoption but minimal peer-reviewed academic treatment, a gap that reflects the broader disconnect between compliance practice and academic research.

The compliance automation research and the Digital Engineering literature thus develop along parallel tracks that have not converged. Compliance researchers build knowledge graphs, property graphs, and domain-specific languages to automate regulatory assessment. Digital Engineering researchers develop model-based approaches, digital threads, and authoritative sources of truth for complex system management. Both communities address documentation accuracy, traceability, and automated verification. Neither community has systematically explored how integrating their approaches might produce capabilities exceeding what either achieves independently. This unexplored intersection is a substantial research opportunity that the present investigation begins to address by establishing whether the professionals who must implement compliance, namely IT and Information Assurance practitioners, recognize value in Digital Engineering capabilities.

A dimension of the compliance imperative that the literature largely leaves implicit is its uneven distribution across organizations. The frameworks above impose the same documentation, traceability, and verification obligations on a community health center, a tribal government, or an under-resourced school district as on a defense prime contractor. Yet the manual labor and specialized expertise those obligations demand are a far heavier proportional burden where staff and budgets are smallest. The compliance burden is, in effect, regressive, and the manual document-centric practices that dominate today entrench that asymmetry. Research commissioned by the United States Small Business Administration documents the pattern across regulation generally: firms with fewer than 20 employees bear an annual federal regulatory cost of \$6,975 per employee, nearly 60\% more than the \$4,463 per employee borne by firms with more than 500 employees~\parencite{Crain_Hopkins_2001}. The pattern extends to cybersecurity compliance specifically. The Government Accountability Office found that concerns about the complexity and cost of CMMC requirements, including third-party assessments estimated at \$4,042 to \$117,768, concentrate among the small businesses that constitute roughly three quarters of the defense industrial base; some of these businesses may stop working with the Department of War rather than absorb those costs~\parencite{GAO_CMMC_2026}.
The few studies that carry model-based methods into resource-limited public-service settings are therefore notable: \citeauthor{Berg_2023} applied Model-Based Systems Engineering to rural-healthcare disaster planning, demonstrating that the approach is not confined to well-funded enterprises and can structure complex requirements in exactly the settings where capacity is scarcest~\parencite{Berg_2023}. If Digital Engineering can substitute reusable model structure for the scarce specialized labor that compliance currently requires, its largest marginal benefit would accrue to small and under-resourced organizations. The equitable distribution of compliance capability thus becomes a substantive, and largely unexamined, stake in whether these methods transfer beyond their sector of origin.

\section{IT Service Management: Structural Failures Despite Mature Frameworks}

IT Service Management frameworks assume capabilities that current practices cannot sustain. ITIL provides comprehensive guidance for aligning IT services with business needs through structured processes, yet its effectiveness depends on the accuracy and currency of underlying configuration information, an accuracy that organizations consistently fail to achieve~\parencite{Cannon_AXELOS_2013}. Version 4 reorganized service management practices and introduced the Service Value System concept. It acknowledges that tracking configurations across virtual systems, cloud computing, and cybersecurity domains presents challenges, yet it provides limited guidance on the documentation accuracy problems that undermine every ITIL process~\parencite{ITIL_4_2019}.

Empirical evidence for ITIL implementation challenges spans two decades of research yet reveals persistent rather than resolving problems. Resistance to change ranks as the top ITIL implementation challenge, at 27\%~\parencite{Cook_2021}. A survey of 491 firms found that while over 90\% use ITSM frameworks, little research examines the benefits actually realized~\parencite{Marrone_Kolbe_2011}. An empirical investigation of ITIL's role in IT governance found that implementation success depends on group efficacy and organizational resources, and that ITIL enables governance processes but ``lacks structural and relational governance practices''~\parencite{Iden_2014}. This finding carries particular weight for the present research: if ITIL provides process governance but lacks structural governance, and if Digital Engineering provides structural governance through model-based system representations, then integrating these approaches addresses a documented deficiency rather than introducing redundant capability.

\citeauthor{Joshi_2022} studied 881 companies worldwide and found that firms with superior IT governance processes achieve more than 25\% higher profits than peer firms with poor governance~\parencite{Joshi_2022}. This finding quantifies the business impact of governance failures, the very failures that ITSM frameworks were designed to prevent but have not resolved. The parallel with MBSE evidence is striking: both MBSE and ITIL are widely adopted based on perceived rather than measured value, and both face adoption challenges rooted in organizational factors rather than technical limitations. Neither literature, however, has put the obvious question to the other: whether the structural governance that model-based representation supplies could repair what ITIL's process governance leaves undone.

Configuration Management Database failures provide the most extensively documented evidence of structural inadequacy, though the headline figures originate from industry-analyst reports rather than peer-reviewed research. That gap itself reflects the practitioner-driven evolution of ITSM documented in the disciplinary history section of this chapter. Gartner reports an 80\% failure rate for CMDB implementations~\parencite{Gartner_CMDB_2019}, and Gartner further predicts that 99\% of organizations using CMDB tooling without addressing data quality gaps will experience visible business disruption~\parencite{Gartner_CMDB_2020}. Forrester research finds that less than half of organizations trust the data in their CMDB~\parencite{Forrester_CMDB_2020}. Peer-reviewed research corroborates these industry findings: \citeauthor{Hauder_Matthes_Roth_2012} found through a study of 123 enterprise architecture practitioners that manual documentation processes cannot maintain accuracy in dynamic environments~\parencite{Hauder_Matthes_Roth_2012}.

Forrester analysis concludes that the CMDB approach itself has failed~\parencite{Forrester_CMDB_Dead_2025}. After decades of implementation attempts across organizations, CMDBs consistently fail to deliver intended value. The analysis attributes failures to structural issues: involving process experts rather than data management professionals, manual data entry that cannot maintain accuracy, and scope creep that renders CMDBs unmanageable. Forrester's research on Application and Infrastructure Dependency Mapping found that 56\% of enterprises report incomplete views of dependencies between applications and underlying infrastructure~\parencite{Forrester_AIDM_2018}. This assessment suggests that incremental CMDB improvements cannot resolve inherent approach limitations. Such a conclusion opens conceptual space for model-based alternatives that maintain accuracy through automated synchronization rather than manual data entry.

Shadow IT compounds the documentation challenge. Gartner research indicates 41\% of employees used shadow IT in 2022, expected to climb to 75\% by 2027, with 30--40\% of large companies' IT expenditure representing shadow IT~\parencite{Gartner_Shadow_IT_2022}. The foundational definition of shadow IT cites industry estimates that its true extent is roughly ten times what CIOs suspect, and identifies key risk factors including security vulnerabilities, compliance violations, and loss of IT control~\parencite{Haag_Eckhardt_2017}.
% Verified 2026-07-02 against the article's first page (author-provided): Haag & Eckhardt 2017 (BISE 59(6):469-473) state "the true extent of shadow IT bypassing the corporate IT units is estimated to be ten times greater than what CIOs suspect within their organizations (Corbin 2015)", a secondhand industry estimate, correctly rendered here as "citing industry estimates."
A systematic review of 126 publications documented how silos between business units and IT departments drive shadow IT, creating security risks, data inconsistency, and compliance violations~\parencite{Klotz_2019}. Systems deployed outside formal governance remain invisible to IT operations and security teams, creating unknown attack surfaces that no documentation process can capture. Examining organizational dynamics when hidden IT systems are discovered reveals interdepartmental conflicts and governance failures that perpetuate documentation gaps~\parencite{Furstenau_2021}. The shadow IT phenomenon reflects a structural mismatch between IT governance and organizational needs: when official IT processes cannot meet user requirements quickly enough, users acquire solutions independently, creating operational dependencies that official documentation does not capture.

Change management effectiveness depends directly on documentation accuracy. Industry analysis confirms that reliance on outdated documentation leads to inaccurate impact assessments~\parencite{Freshservice_Change_2021}. Research by \citeauthor{Bokan_Santos_2021} highlights difficulties organizations encounter in maintaining comprehensive security oversight~\parencite{Bokan_Santos_2021}. Change management depends on accurate understanding of system relationships~\parencite{iso20000_2018} that current documentation approaches cannot maintain. Each change approved against inaccurate documentation carries incident risk. When impact assessments miss dependencies that exist in operational systems, changes cause unintended effects that consume resources, damage trust in change processes, and create pressure for emergency changes that further degrade documentation accuracy. The resulting cycle is self-reinforcing and resistant to process improvement within the document-centric paradigm.

Recent applications confirm that Model-Based Systems Engineering continues to expand into new and demanding complex-systems problems. \citeauthor{SUN2025100231}, for instance, developed an interactive visualization-based MBSE technique for requirements analysis and test design in the joint testing of complex systems, reporting gains in the accuracy and dependability of the resulting test plans~\parencite{SUN2025100231}. Within the enterprise domain specifically, prior research by \citeauthor{Bonar_Hastings_2024} presented a preliminary reference model in which Digital Engineering practices support compliance verification~\parencite{Bonar_Hastings_2024}. Yet no peer-reviewed study has directly examined integrating Model-Based Systems Engineering with IT Service Management. That absence is consequential; it is the precise gap this dissertation confronts.

\section{Enterprise Visibility: Empirical Evidence of Systemic Failure}\label{sec:enterprise_visibility}

The challenges documented in the preceding sections share common roots in the inability of organizations to maintain accurate, current documentation of their enterprise information systems. This section synthesizes peer-reviewed research and industry analysis to establish that these challenges reflect systemic patterns rather than isolated organizational deficiencies, and to provide the empirical foundation for evaluating Digital Engineering as a potential disciplinary response.

Independent studies converge on a single finding: organizations cannot see large portions of their own IT environments, and the blind spots concentrate in the cloud. Monitoring coverage is incomplete at the outset. One global study found that organizations can monitor only 66\% of their IT environments, leaving blind spots particularly in cloud deployments~\parencite{IDC_Exabeam_2023}, and 63\% of security teams lack visibility and control into all user device activity connected to their infrastructure~\parencite{Ponemon_Visibility_2023}. Practitioner confidence tracks that shortfall. Only 15\% of respondents to the SANS Institute SOC Survey expressed very high confidence that all devices on their network are discoverable~\parencite{SANS_SOC_Survey_2023}. The gaps then compound across organizational boundaries, because 55\% of organizations maintain security and IT data silos and 62\% report that those silos slow security response times~\parencite{Ivanti_Cybersecurity_2025}. Cloud infrastructure concentrates the exposure most sharply. Cloud-related breaches struck 95\% of organizations in the preceding eighteen months~\parencite{CSA_Cloud_Security_2024}, 82\% of enterprises experienced security incidents due to cloud misconfigurations, and 67\% struggle with limited visibility into cloud infrastructure~\parencite{Check_Point_Cloud_2024}. The scale of the unseen surface is itself unknown. Organizations use 15--22 times more cloud services than their IT departments estimate~\parencite{Cisco_Cloud_Services_2016}, and Gartner predicted that through 2025, 99\% of cloud security failures would be the customer's fault~\parencite{Gartner_Cloud_Misconfiguration_2024}.
% Verified: Gartner's prediction reads "Through 2025, 99% of cloud security failures will be the customer's fault," from "Is the Cloud Secure?" (Smarter with Gartner, 2019); confirmed via CIO.com (cio.com/article/416343) quoting and linking the Gartner piece (gartner.com blocks automated access). Horizon corrected from 2027 to 2025; bib entry Gartner_Cloud_Misconfiguration_2024 needs updating to the 2019 Gartner source.

Peer-reviewed research provides empirical grounding for these industry findings. An empirical study of configuration errors in commercial and open source systems found that 70--85\% of misconfigurations result from mistakes in setting configuration parameters and that 22--57\% involve configurations external to the examined system~\parencite{Yin_Yuan_Lu_2011}. NIST Special Publication 800-128 defines configuration drift as systems deviating from baseline configurations over time through manual interventions, software updates, and environmental factors~\parencite{NIST_SP_800_128}. The Uptime Institute's Annual Outage Analysis confirms that 64\% of IT system outages occurred because of configuration or change management issues~\parencite{Uptime_Outage_2023}. The IT Process Institute established that 80\% of unplanned outages result from ill-planned changes made by administrators or developers~\parencite{IT_Process_Institute_2004}. Proper dependency documentation would have flagged many such changes.

Breach detection times serve as proxy measures for organizational visibility. \citeauthor{IBM_Ponemon_2024} report that the mean time to identify breaches reached 204 days, with breaches involving lifecycles exceeding 200 days costing an average of \$5.46 million~\parencite{IBM_Ponemon_2024}. The 2025 report found that 35\% of breaches involved shadow data, defined as information stored in unmanaged locations, and 40\% involved data stored across multiple environments that organizations struggle to inventory comprehensively~\parencite{IBM_Ponemon_2025}. These findings demonstrate that visibility failures directly impact security outcomes in measurable economic terms.

Theoretical and empirical frameworks help explain why these failures persist, and the enterprise architecture literature shows the same failure mechanisms recurring across independent studies. Complexity-driven degradation is the first mechanism: through PLS-SEM analysis of 249 information systems managers, \citeauthor{Beese_2023} demonstrated that IS architecture complexity significantly reduces efficiency, flexibility, transparency, and predictability, and that organizations without adequate enterprise architecture management face compounding architectural degradation~\parencite{Beese_2023}. Unvalidated framework prescriptions constitute a second mechanism: an examination of twenty-seven organizations by \citeauthor{Kotusev_2019} found that prescribed enterprise architecture artifact lists from frameworks like TOGAF were never empirically validated and do not reflect actual practice, documenting a fundamental gap between what frameworks recommend and what organizations find useful~\parencite{Kotusev_2019}. Stakeholder-engagement inhibitors form a third: \citeauthor{Kurnia_2021} identified organizational silos, political barriers, and poor cross-team communication as inhibitors that directly undermine documentation quality and IT governance effectiveness~\parencite{Kurnia_2021}. In a scoping review of forty studies on organizational silos, \citeauthor{Bento_Tagliabue_Lorenzo_2020} identified five conceptualizations and characterized silo mentality as an absence of systems thinking and organizational vision~\parencite{Bento_Tagliabue_Lorenzo_2020}. Documentation challenges cut across all three mechanisms: \citeauthor{Bree_Karger_2022} conducted a systematic literature review organizing enterprise architecture management challenges into six dimensions, with documentation challenges including dearth of automated tools, immature documentation models, and insufficient emphasis on forward-looking documentation~\parencite{Bree_Karger_2022}.

A complementary stream examines the pathways through which enterprise architecture delivers value when it succeeds, and it reaches equally sobering conclusions. \citeauthor{Tamm_2022}, drawing on twenty-two expert interviews and two in-depth case studies, investigated how enterprise architecture leads to organizational benefits, observing that the mechanisms underpinning EA benefit realization remain ``not fully understood''~\parencite{Tamm_2022}. \citeauthor{van_de_Wetering_2022}---in motivating a survey of 414 senior practitioners, business and IT managers, and executives---observes that despite high expectations ``many EA initiatives in practice fail or do not live up to stakeholder expectations,'' reinforcing the pattern of enterprise architecture as a discipline with acknowledged value but unreliable execution~\parencite{van_de_Wetering_2022}.
% Verified: "not fully understood" is Tamm et al.'s wording (abstract, doi:10.1016/j.ijinfomgt.2022.102554: "the mechanisms underpinning EA benefit realisation are not fully understood"; the twenty-two interviews and two in-depth case studies are also confirmed there). The "many EA initiatives in practice fail or do not live up to stakeholder expectations" sentence is van de Wetering's own (Heliyon 8(11):e11484, full text at pmc.ncbi.nlm.nih.gov/articles/PMC9640386), citing Al-Kharusi et al. 2021, Ross et al. 2019, and Tormer and Henningsson 2019, not Tamm; attribution corrected. n=414 confirmed; van de Wetering motivates rather than empirically documents the shortfall, so "documenting widespread shortfalls" was softened.
A fuzzy-set qualitative comparative analysis investigated how enterprise architecture management contributes to digital transformation capabilities, identifying configurational paths to effectiveness that require combinations of factors rather than any single intervention~\parencite{Pattij_2022}.

Complexity theory provides an explanatory framework for why traditional documentation approaches fail in modern enterprise environments. Complex Adaptive Systems theory, applied to information systems development, yields seven first principles of adaptive success and the argument that if complexity is not managed appropriately, information systems fail~\parencite{Benbya_McKelvey_2006}. Later work demonstrated that enterprise information systems have reached complexity levels exceeding prior technological generations~\parencite{Benbya_Nan_Tanriverdi_Yoo_2020}. Enterprise IT environments exhibit characteristics of complex adaptive systems, including many interconnected components, emergent behaviors arising from component interactions, continuous change, and unpredictable responses to interventions. Static documentation approaches assume that these characteristics remain stable between documentation updates.

Technical debt research adds a final perspective, and across primary studies, mapping studies, and secondary syntheses it converges on treating documentation debt as a distinct, recurring category rather than an incidental concern. One study specifically addressed documentation technical debt, defined as problems concerning non-existent, inadequate, or incomplete software project documentation~\parencite{Santos_2019}. A systematic mapping study by \citeauthor{Li_Avgeriou_Liang_2015}, covering ninety-four studies, established technical debt taxonomy and management frameworks that include documentation debt as a distinct category~\parencite{Li_Avgeriou_Liang_2015}. Across nineteen secondary studies, \citeauthor{Junior_Travassos_2022} consolidated perspectives on technical debt~\parencite{Junior_Travassos_2022}, and \citeauthor{Kleinwaks_2023} extended technical debt concepts to systems engineering contexts~\parencite{Kleinwaks_2023}. Documentation technical debt accumulates when organizations defer documentation updates to prioritize operational activities. Unlike code technical debt, documentation debt remains invisible until documentation is needed for change impact assessment, incident response, or compliance verification. At that point, the accumulated debt imposes costs far exceeding the original deferral savings.

Table~\ref{tab:visibilityStatistics} summarizes the evidence documenting enterprise visibility and documentation failures across peer-reviewed and industry research sources.

\begin{table}[htbp]
  \centering
  \caption{Enterprise visibility and documentation failure evidence.}\label{tab:visibilityStatistics}
  % ALT: Table of enterprise visibility and documentation failure findings grouped into visibility-gap, configuration-management, shadow-IT, and security-impact categories; each row gives a finding, its statistic, and the source.
  \small
  \begin{tabular}{@{}p{3.0in}p{1.0in}p{2.0in}@{}}
    \toprule
    \textbf{Finding} & \textbf{Statistic} & \textbf{Source} \\
    \midrule
    \multicolumn{3}{@{}c@{}}{\textit{Visibility Gap Metrics}} \\
    IT environment monitorable & 66\% & \parencite{IDC_Exabeam_2023} \\
    Security teams lacking device visibility & 63\% & \parencite{Ponemon_Visibility_2023} \\
    High confidence in device discovery & 15\% & \parencite{SANS_SOC_Survey_2023} \\
    Organizations with security/IT silos & 55\% & \parencite{Ivanti_Cybersecurity_2025} \\
    \addlinespace
    \multicolumn{3}{@{}c@{}}{\textit{Configuration Management Failures}} \\
    CMDB implementation failure rate & 80\% & \parencite{Gartner_CMDB_2019,Gartner_CMDB_2020} \\
    Outages from configuration issues & 64\% & \parencite{Uptime_Outage_2023} \\
    Misconfigurations from parameter errors & 70--85\% & \parencite{Yin_Yuan_Lu_2011} \\
    Unplanned outages from ill-planned changes & 80\% & \parencite{IT_Process_Institute_2004} \\
    \addlinespace
    \multicolumn{3}{@{}c@{}}{\textit{Shadow IT and Undocumented Assets}} \\
    Shadow IT as percentage of IT spend & 30--40\% & \parencite{Gartner_Shadow_IT_2022} \\
    Cisco cloud services vs. IT estimates & 15--22x higher & \parencite{Cisco_Cloud_Services_2016} \\
    Employees using shadow IT (2022) & 41\% & \parencite{Gartner_Shadow_IT_2022} \\
    Projected shadow IT usage (2027) & 75\% & \parencite{Gartner_Shadow_IT_2022} \\
    \addlinespace
    \multicolumn{3}{@{}c@{}}{\textit{Security Impact Metrics}} \\
    Mean time to identify breach & 204 days & \parencite{IBM_Ponemon_2024} \\
    Cloud breaches from misconfigurations & 82\% & \parencite{Check_Point_Cloud_2024} \\
    Organizations with cloud breaches (18 mo) & 95\% & \parencite{CSA_Cloud_Security_2024} \\
    Cloud security failures attributed to customer fault (through 2025) & 99\% & \parencite{Gartner_Cloud_Misconfiguration_2024} \\
    \bottomrule
  \end{tabular}

  \smallskip
  {\footnotesize \textit{Note.} The 80\% CMDB implementation failure rate derives from industry-analyst (Gartner) reports rather than peer-reviewed research and should be read as practitioner-oriented estimation rather than empirically validated measurement.}
\end{table}

Peer-reviewed empirical research and industry analysis converge to establish that enterprise visibility and documentation failures represent a systemic challenge rather than isolated organizational deficiencies. Traditional documentation approaches---including manual configuration tracking, periodic documentation updates, and static architecture diagrams---cannot maintain accuracy in environments characterized by continuous change, complex dependencies, and organizational silos.

The evidentiary sources in this section warrant a methodological note. The enterprise visibility evidence draws substantially on industry analyst reports and practitioner surveys alongside peer-reviewed academic research. This reliance on industry gray literature reflects a characteristic of the research domain rather than an analytical preference: the operational challenges of enterprise IT documentation and visibility have received considerably more attention from industry analysts and practitioner communities than from academic researchers. The relative scarcity of peer-reviewed empirical studies quantifying CMDB failure rates, shadow IT prevalence, and breach detection timelines itself constitutes evidence of the research gap this dissertation addresses. Peer-reviewed evidence corroborates industry findings in several instances, as with the work of \citeauthor{Hauder_Matthes_Roth_2012} on documentation accuracy, \citeauthor{Santos_2019} on documentation technical debt, \citeauthor{Beese_2023} on architecture complexity, \citeauthor{Tamm_2022} on enterprise architecture failures, and \citeauthor{Kotusev_2019} on enterprise architecture practice. Where such corroboration occurs, the convergence strengthens confidence in the broader pattern. This review presents the industry evidence not as a substitute for academic rigor but as the best available evidence for phenomena that academic research has not yet systematically investigated.

\section{Research Gaps and Theoretical Framework}

The structured review of the literature reveals a pattern: frameworks and compliance requirements assume documentation and visibility capabilities that organizations frequently lack, while Digital Engineering methods that could address these capabilities remain confined within a discipline that enterprise IT and Information Assurance have not engaged. This section synthesizes the findings into a theoretical framework while documenting the research gaps that this investigation begins to address.

The academic research gap is pronounced. Systematic literature reviews examining MBSE report no published research, within the scope of their searches, addressing enterprise IT infrastructure or Information Assurance applications as an integrated enterprise discipline, as distinct from MBSE applied to individual engineered systems or to single compliance controls, beyond the preliminary reference model by \textcite{Bonar_Hastings_2024} and the conceptual DevSecOps framework by \textcite{CMU_SEI_DevSecOps_2023}. This gap persists despite explicit requirements in compliance frameworks for capabilities that Digital Engineering provides. As the disciplinary history presented in this chapter demonstrates, the gap cannot be attributed to Digital Engineering immaturity, given that defense and aerospace have employed Digital Engineering successfully for years. Nor can it be attributed to tool unavailability, since MBSE tools have existed for over a decade and PLM systems for several decades, with digital twin platforms maturing rapidly over the past ten years. Rather, the gap reflects disciplinary boundaries that formed when systems engineering, Information Assurance, cybersecurity, and IT service management each developed independent academic identities, professional communities, and research traditions. The low interconnectedness between organizational units that Henderson and Salado associate with impeded adoption operates at the disciplinary level itself~\parencite{Henderson_Salado_2024_Org}.

Standards bodies have recognized enterprise applicability of systems engineering approaches. The Unified Architecture Framework provides viewpoints applicable to enterprise IT\@. NIST publications require enterprise architecture capabilities for compliance. ITIL requires visibility and documentation that model-based approaches could provide. Yet academic research has not examined practical application. This standards-research disconnect leaves practitioners without empirical guidance for applying available standards to enterprise IT challenges.

Table~\ref{tab:researchGaps} summarizes the research gaps identified across the literature domains examined in this review.

\begin{table}[htbp]
  \centering
  \caption{Research gaps within corpus of knowledge.}\label{tab:researchGaps}
  % ALT: Table of research gaps by literature domain; each row gives the domain, a description of the gap, and its research implication.
  \small
  \begin{tabular}{@{}p{1.8in}p{2.1in}p{2.2in}@{}}
    \toprule
    \textbf{Domain} & \textbf{Gap Description} & \textbf{Research Implication} \\
    \midrule
    MBSE for Enterprise IT & One study applying MBSE to enterprise IT, none for IA program management & Foundation research required \\
    MBSE Adoption Evidence & Adoption studies limited to SE populations; no data from IT or IA professionals & Perception measurement needed in new populations \\
    Disciplinary Integration & No research connecting SE, IA, cybersecurity, and ITSM methods despite shared challenges & Cross-disciplinary investigation needed \\
    Digital Threads & No research on traceability for IT/IA contexts & Conceptual validation needed \\
    Digital Twin & Growing security application research but no enterprise IT integration studies & Application studies needed \\
    ITSM Integration & No frameworks integrating DE with ITIL & Integration research required \\
    Compliance Automation & Model-based compliance research and DE research develop on parallel tracks without convergence & Integration research needed \\
    Adoption Theory & Cybersecurity adoption and MBSE adoption research use same theories but never intersect & Bridging research needed \\
    Open Source & No academic validation for enterprise IT & Evaluation research needed \\
    Professional Perceptions & Unknown awareness and perceived value among IT/IA professionals & This research addresses \\
    \bottomrule
  \end{tabular}
\end{table}

Based on the structured review of the literature, this research adopts a theoretical framework integrating Digital Engineering principles with established Information Assurance and IT Service Management practices. This framework posits that Digital Engineering is a disciplinary approach with capabilities reported in defense and aerospace contexts that align with gaps that have persisted in enterprise IT despite decades of framework development and organizational investment. The framework does not assert that Digital Engineering will resolve these gaps in enterprise IT contexts, which remains an empirical question. Consistent with the boundary-object and knowledge-transformation lens established earlier in this chapter~\parencite{Star_Griesemer_1989, Carlile_2004}, the framework treats the structural correspondences developed below not as mere analogies but as the conditions under which Digital Engineering models could function as boundary objects. Such boundary objects would span the pragmatic knowledge boundary~\parencite{Carlile_2002, Carlile_2004} that separates systems engineering from Information Assurance and IT service management. It thereby identifies structural correspondences between demonstrated Digital Engineering capabilities and documented enterprise IT challenges, establishing the theoretical basis for investigating whether practitioners recognize them.

Persistent failures documented in IT Service Management and Information Assurance practices share common root causes that Digital Engineering practices are designed to address. Organizations struggle with documentation accuracy because traditional approaches rely on manual processes disconnected from operational systems. Organizations fail to maintain traceability because document-centric methods cannot sustain verified connections as systems evolve. Organizations lack visibility because static artifacts cannot represent dynamic system states. Digital Engineering's reported ability, measured in rare but consistent cases, to address these root causes in defense and aerospace contexts motivates investigation of whether similar benefits may transfer to enterprise IT environments.

The Department of Defense's 2018 Digital Engineering Strategy defines the authoritative source of truth as ``a single source of data and models'' that provides ``a definitive technical baseline'' for programs~\parencite{DoD_DE_Strategy_2018}. Current IT and Information Assurance practices lack authoritative sources of truth, instead maintaining multiple disconnected documentation artifacts that diverge over time. Digital threads establish and maintain authoritative traceability throughout system lifecycles, addressing the gap between security requirements, control implementations, and compliance evidence. Digital twin capabilities enable organizations to simulate system behavior, test proposed changes, and analyze scenarios without affecting production systems. Model-based approaches maintain synchronization with operational systems, providing the visibility that manual documentation cannot sustain.

The theoretical framework acknowledges several limitations. First, it extrapolates from Digital Engineering value demonstrated in aerospace and defense to anticipated value in enterprise IT contexts. Whether benefits demonstrated for physical systems transfer to logical information systems remains unvalidated. Second, it assumes that Digital Engineering tools and methods can be adapted for enterprise IT contexts. \citeauthor{Mazeika_Butleris_2020} documented that current MBSE methods inadequately address security requirements~\parencite{Mazeika_Butleris_2020}, suggesting adaptation requirements may exceed anticipated effort. Third, the framework does not address organizational change management, workforce development, or cultural transformation requirements. The adoption barriers documented by \citeauthor{Henderson_McDermott_Salado_2024}---including middle management resistance, awareness deficits, and the need for dedicated MBSE teams~\parencite{Henderson_McDermott_Salado_2024}---would likely manifest with greater intensity in domains unfamiliar with systems engineering practices. Fourth, the framework focuses on potential benefits without comprehensive analysis of costs or implementation challenges. Cost-benefit analysis requires empirical data this research does not collect.

\section{Chapter Summary}

This chapter has examined the literature across interconnected domains relevant to applying Digital Engineering methods to Information Assurance and IT Service Management, advancing four interrelated arguments that position the present research within the academic conversation.

First, the independent evolution of systems engineering, Information Assurance, cybersecurity, and IT service management as separate academic and professional disciplines created structural barriers to cross-disciplinary methodological exchange. Systems engineering emerged from defense laboratories and engineering departments; Information Assurance grew from national security policy and computing education; cybersecurity crystallized as a meta-discipline drawing from multiple computing fields; and IT service management developed as a practitioner-driven domain with limited academic theorization. These parallel trajectories produced separate journals, conferences, curricular standards, and professional communities, institutional structures that explain the absence of research connecting Digital Engineering to enterprise IT and Information Assurance.

Second, the evidence base for MBSE value, while compelling in individual cases such as the 18\% efficiency improvement and 9\% interface defect reduction documented by \textcite{Rogers_Mitchell_2021}, remains disproportionately reliant on perceived rather than measured outcomes. \citeauthor{Henderson_Salado_2021} found that approximately two-thirds of claimed benefits lack empirical measurement~\parencite{Henderson_Salado_2021}. Efforts to develop measurement frameworks by \citeauthor{Henderson_2023_Metrics} and evaluations of perceived value by \citeauthor{Campo_2023} have advanced the methodological conversation without resolving the measurement deficit~\parencite{Henderson_2023_Metrics, Campo_2023}. This evidence deficit creates rational hesitation among organizations outside the systems engineering discipline, where institutional experience cannot compensate for the measurement gap.

Third, adoption research demonstrates that perceptions govern adoption decisions more powerfully than objective technical merit. \citeauthor{Call_2024} showed that perceived complexity and compatibility barriers impede adoption even among systems engineering professionals who recognize MBSE's relative advantage~\parencite{Call_2024}. \citeauthor{Henderson_McDermott_Salado_2024} found that approximately 22\% of interviewed systems engineering professionals could not clearly define MBSE~\parencite{Henderson_McDermott_Salado_2024}. Parallel research in cybersecurity contexts by \citeauthor{Hasani_2023} and \citeauthor{Anthony_2024} demonstrates that the same adoption theories and barrier patterns operate across technology domains~\parencite{Hasani_2023, Anthony_2024}. If awareness deficits and perception barriers persist within the originating discipline, their magnitude among IT and Information Assurance professionals, who operate outside the systems engineering discipline entirely, demands empirical measurement rather than assumption.

Fourth, compliance frameworks and IT service management standards require capabilities that current practices consistently fail to provide, while Digital Engineering offers those capabilities in adjacent domains. The enterprise visibility evidence documents the scope of failure: 66\% IT environment monitoring, the industry-reported 80\% CMDB failure rate, and 204-day average breach detection times. Enterprise architecture research by \citeauthor{Tamm_2022} and \citeauthor{van_de_Wetering_2022} confirms that EA initiatives frequently fail to meet stakeholder expectations~\parencite{Tamm_2022, van_de_Wetering_2022}. The compliance automation literature, including emerging work on continuous compliance by \citeauthor{Santilli_2023} and \citeauthor{Angermeir_2024}, demonstrates growing recognition that model-based approaches can address documentation and traceability requirements~\parencite{Santilli_2023, Angermeir_2024}. Yet no research connects these parallel tracks by examining Digital Engineering as an integrated solution for enterprise IT and Information Assurance challenges.

The theoretical framework developed from this synthesis posits that structural correspondences between demonstrated Digital Engineering capabilities and documented enterprise IT challenges warrant investigating whether practitioners recognize these correspondences. The research contributes by investigating whether IT and Information Assurance professionals perceive value in Digital Engineering capabilities. If professionals perceive value, findings justify subsequent implementation research. If professionals do not perceive value despite documented challenges, findings challenge the theoretical premise and identify barriers the field must address before adoption can occur.

Chapter 3 presents the research methodology employed to investigate professional awareness and perceptions of Digital Engineering capabilities. The methodology uses a quantitative survey-based approach following a systems engineering lifecycle to ensure rigor and traceability throughout the research process.